% s-298-p298-m002; original PDF 298, 299, 300, 301
\Aphorism{s-298-p298-m002}{356}
Vielleicht ist noch nie die philosophische Wahrheit despotischer auf den Kopf gestellt worden als mit dem Satze, den Kant in der „Grundlegung zur Metaphysik der Sitten“ gewagt hat: „Empirische Principien taugen überall nicht dazu, um moralische Gesetze darauf zu gründen.“ Wir wollen uns hier nicht mit Kants Moral auseinandersetzen; vermuthlich ist gegen den kategorischen Imperativ noch manches Andere einzuwenden als seine metaphysische Herleitung, die uns für jetzt belustigen soll. Denn dass es daran etwas zum Lachen giebt, wird auch der ernsthafteste Moralpredigtamts-Candidat zugestehen müssen; man male sich nur die Scene aus! Gesucht ein kategorischer Imperativ; aber kategorisch muss er sein, denn hypothetische Imperative taugen überall nicht dazu u. s. w. Moral muss \Sperr{unbedingt} befehlen, sonst ist es keine Moral für Ostpreussen. Erster Versuch: die Sinnenwelt. Aber da giebt es zwar Imperative, jedoch keine kategorischen; da herrscht Heteronomie des Willens, der Quell aller unächten Principien der Sittlichkeit. Bleibt nur die intelligible Welt (d. h. diejenige, die nach der Kritik der reinen Vernunft \Sperr{nicht intelligibel} ist), und nun lässt sich zwar das Kategorische retten, mit einer leichten quaternio terminorum zwischen (begrifflicher) Allgemeinheit eines Gesetzes und (socialer) Allgemeinheit seines Anwendungsbereichs, aber — der Imperativ ist dabei verloren gegangen! Oder glaubte Kant im Ernst an die Möglichkeit eines „Du sollst“, abzuleiten und zu rechtfertigen aus dem Dasein einer Welt, die er selbst nicht etwa nur als quantitativ unendlich, sondern als qualitativ transcendent begriffen hatte? Schon die pluralité des mondes im astronomischen, also immanenten Sinne entmuthigt ja ein feineres Gewissen und verhöhnt die grossartige Würde und Selbstgefälligkeit, mit der wir unsere tellurische Moralwirthschaft betreiben; ja, wie es scheint, fehlt es uns noch innerhalb viel engerer Grenzen an einigen Vorbedingungen zur Berechtigung, den objectiven Werth unserer Handlungen abzuschätzen oder gar sehr hoch abzuschätzen. In wie weit kommt, mit allgemeineren Massstäben gemessen, der Einzelne überhaupt in Betracht? Ist er im Stande, der Art merklich zu nützen oder zu schaden? Sind böse und gute Handlungen vielleicht zum Bestande des \Sperr{Lebens} gleich nothwendig? und so gäbe es noch einige andere Fragen, vor deren Beantwortung man keinesfalls behaupten darf, dass die individuelle Willenshandlung auch nur im Irdischen sonderlich ins Gewicht falle. Indessen, vielleicht belehrt uns hier die fernere Erfahrung eines Anderen; und was die kosmische Unendlichkeit des Daseins anbelangt, so reducirt sie den objectiven Werth unserer menschlichen Thaten und Unthaten zwar \Sperr{beinahe} auf Null, aber doch nur beinahe, und sie rettet ihm wenigstens seine quantitative Vergleichbarkeit mit dem Totalwerthe des Ganzen. Was aber sagt die transcendente Welt zu den Voraussetzungen der Moral? Jene Welt, die nach Kant „Etwas \Math{= X}“, d. h. vollständig unerkennbar ist und sich demnach auch in Folge unserer Handlungen in völlig unerkennbarer, uncontrollirbarer Richtung verändert, also empirisch gerechnet \Sperr{gar nicht} verändert? Wäre diese Welt, deren Gefühlsabbild nur der radicale Indifferentismus von der menschenunmöglichsten Härte und Jenseitigkeit sein könnte, wirklich zum Ankergrund einer moralischen Auswerthung des Daseins geeignet, einer Auswerthung, die uns aufs entschiedenste an das Gegentheil des Indifferentismus, an einen höchst unterschiedlichen, höchst ungleichgültigen, höchst wesentlichen und definitiven Sinn unseres Wollens und Handelns glauben machen will? Ist es nicht grausam, den Menschen zur Formung und Gestaltung der Dinge anzutreiben und ihm eine amorphe Gallertmasse in die Hand zu geben, die keine oder jede Form annimmt? ist es nicht logischer Radbruch und Hexen-Einmaleins, aus der Erkenntniss, dass kein menschliches Wissen und Wirken den Kern der Welt trifft, zu folgern: handle moralisch, \Sperr{denn} deine Handlung hat einen Werth „an sich“? Ist diese Kantische Moral nicht ein Nachschössling jener Prädestinationslehre, die dem Menschen in Einem Athem sagt „über Heil und Verdammniss deiner Seele ist längst entschieden“ und zugleich „von deinem gegenwärtigen Thun und Glauben hängt Heil oder Verdammniss ab“? Erkennen wir dieses christliche \Foreign{credo quia absurdum} nicht allzugut wieder in jenem Kantischen \Foreign{nego quia non absurdum,} zu deutsch: empirische Principien taugen überall nicht dazu, um moralische Gesetze darauf zu gründen!
\par

